\corrigesChapitre{9}{Lois discrètes et loi binomiale}
\begin{corrige}{1}{Uniforme discrète}
\begin{questions}
\item $X$ suit la loi uniforme sur $\{1,2,3,4,5,6\}$ ; chaque probabilité vaut $1/6$.
\item $(6+1)/2=3{,}5$.
\item $2/6=1/3$, car les issues $5$ et $6$ conviennent.
\end{questions}
\end{corrige}
\begin{corrige}{2}{Bernoulli}
\begin{questions}
\item Bernoulli de paramètre $0{,}3$.
\item $E(X)=0{,}3$ et $V(X)=0{,}3\times0{,}7=0{,}21$.
\item $\sqrt{0{,}21}\simeq0{,}458$.
\end{questions}
\end{corrige}
\begin{corrige}{3}{Pascal}
\begin{questions}
\item $1,4,6,4,1$.
\item $\binom52=4+6=10$.
\item $\binom53=\binom52=10$.
\end{questions}
\end{corrige}
\begin{corrige}{4}{Identifier le modèle}
\begin{questions}
\item $\mathcal B(10,1/2)$, car les répétitions sont indépendantes et de même paramètre.
\item La composition de l’urne change, donc les probabilités conditionnelles changent et l’indépendance manque.
\item $n=20$ : il compte les répétitions, pas la probabilité du succès.
\end{questions}
\end{corrige}
\begin{corrige}{5}{Une loi complète}
\begin{questions}
\item Pour $k=0,1,2,3$, les probabilités sont $1/8,3/8,3/8,1/8$.
\item $3/8+1/8=1/2$.
\item $0+3/8+6/8+3/8=12/8=3/2=np$.
\end{questions}
\end{corrige}
\begin{corrige}{6}{Au moins un défaut}
\begin{questions}
\item $X\sim\mathcal B(20,0{,}02)$.
\item $1-0{,}98^{20}\simeq0{,}332392$.
\item $E(X)=0{,}4$ défaut par lot en moyenne théorique sur de nombreux lots ; ce n’est pas un nombre de défauts réalisable dans un lot.
\end{questions}
\end{corrige}
\begin{corrige}{7}{Cumuler}
\begin{questions}
\item $0{,}8^4=0{,}4096$ et $4(0{,}2)(0{,}8)^3=0{,}4096$.
\item $0{,}8192$.
\item $1-0{,}8192=0{,}1808$.
\end{questions}
\end{corrige}
\begin{corrige}{8}{Choisir une taille}
\begin{questions}
\item $1-0{,}9^n$.
\item $0{,}9^n\leq0{,}05$, donc $n\geq\ln(0{,}05)/\ln(0{,}9)\simeq28{,}433$ ; il faut $29$ tentatives.
\item Les succès peuvent se produire ensemble ; les événements ne sont pas incompatibles. L’addition naïve donnerait même des valeurs supérieures à $1$.
\end{questions}
\end{corrige}
\begin{corrige}{9}{Gagner ou perdre}
\begin{questions}
\item $G=5X-2$ : gain $3$ en cas de succès, perte $2$ sinon.
\item $E(G)=5(0{,}3)-2=-0{,}5$ euro.
\item Le nombre de réussites $S$ suit $\mathcal B(100,0{,}3)$ ; $E(5S-200)=150-200=-50$ euros.
\end{questions}
\end{corrige}
\begin{corrige}{10}{Prouver la symétrie}
\begin{questions}
\item Les places d’échec sont exactement le complément des places de succès. Le choix est bijectif.
\item Les deux ensembles de choix ont donc le même nombre d’éléments.
\item Non : les facteurs $p^k(1-p)^{n-k}$ et $p^{n-k}(1-p)^k$ diffèrent généralement. La symétrie des masses est assurée pour $p=1/2$.
\end{questions}
\end{corrige}
\begin{corrige}{11}{Un paramètre inconnu}
\begin{questions}
\item $12p=3$, donc $p=1/4$.
\item $\sigma=\sqrt{12(1/4)(3/4)}=3/2$.
\item $P(X=0)=(3/4)^{12}\simeq0{,}031676$.
\end{questions}
\end{corrige}
\begin{corrige}{12}{Détecter un faux modèle}
\begin{questions}
\item Les tirages ne sont pas indépendants ; après un tirage correct, la proportion de défauts change.
\item $(8/10)(7/9)(6/8)=7/15\simeq0{,}466667$.
\item $0{,}8^3=0{,}512$ est différent. Une approximation ne doit pas être présentée comme un calcul exact sans justification.
\end{questions}
\end{corrige}
